\section{Introduction}
\label{sec:intro}

In a partially observable environment the current observation can be identical in situations that call for different actions. This is perceptual aliasing \citep{whitehead1991aliasing,mccallum1996thesis}. An agent must then act on a summary of its history. Predictive state representations make the point precise: a history is summarised by predictions of \emph{action-conditional} tests, that is, of future observations given future actions \citep{littman2001psr,singh2004psr}. Neural descendants of this idea train recurrent history encoders with multi-step, action-conditional prediction losses \citep{gregor2019simcore,guo2018belief,guo2020pbl,kwon2024future}. Related latent-consistency objectives with action-conditional dynamics are standard in model-based reinforcement learning \citep{schwarzer2021spr,hansen2024tdmpc2,schrittwieser2020muzero}.

An action-conditional target is only observed for the actions that were actually taken. With a resettable simulator one can do more: save the complete simulator state, try several first actions from it, and restore the state between attempts. This \emph{branching} yields paired counterfactual targets from the same history. It is an old device, used for instance by the vine procedure of TRPO \citep{schulman2015trpo} and by PSR learning with resets \citep{james2004reset}. Branching is not free, however. Each branch step is an environment transition, and the restore operation is a privilege that ordinary interaction with a real system does not provide. Comparisons of representation learners that ignore either cost can credit an objective with gains that actually come from extra experience or from simulator access.

This draft asks a narrow question: \emph{under an equal environment-transition budget and a fixed readout protocol, how do the effects of branching collection and of action conditioning separate?} We do not propose a new objective. Action-conditional prediction, multi-step prediction, recurrent world models, and state-restore branching are all prior work (Section~\ref{sec:related}). Our contributions are limited to the following.
\begin{itemize}
  \item \textbf{An accounting protocol and implementation.} A single transition ledger charges every main-trajectory, branch, and evaluation step against one hard cap. Snapshots and restores are counted separately and are forbidden during evaluation. Actors and learners receive only the observable history, never simulator state or evaluator labels. These properties are checked by unit tests. The tests are technical checks, not evidence about any method.
  \item \textbf{An evaluator diagnostic.} Before comparing learned encoders, we check that the readout (a frozen encoder, a fresh action-conditional reward head, and an exhaustive model-predictive planner) can separate informative from uninformative representations. It can: a hand-designed history-only control and an untrained recurrent encoder differ by \DiagHeadroomMin{} in forced-commit accuracy in every diagnostic seed (Section~\ref{sec:diag}).
  \item \textbf{A small controlled 2$\times$2 study} that crosses action conditioning with branching under an equal transition budget and reports experience and compute separately (Section~\ref{sec:study}).
\end{itemize}
We also report what the evidence does not show. The primary metric, forced-commit accuracy, measures a single junction decision, not planning success in general. The full-action planner's success rate is confounded by stalling behaviour. All results come from one synthetic task family with three seeds, and the comparison with an external action-conditioned method such as that of \citet{kwon2024future} is not yet run.
